\chapter{Quotient Filters and Expansion}

The previous chapter established that Bloom filters cannot delete and 
cuckoo filters cannot expand. Both limitations stem from the same root 
cause: the stored representation (bits or fingerprints) does not retain 
enough information to reconstruct the element's position in a resized 
table.

Quotient filters, introduced by Bender et al.\ \cite{quotient}, solve 
the expansion problem with an elegant insight: instead of computing the 
bucket index and fingerprint from \emph{separate} parts of the hash 
(as cuckoo filters do), they derive both from a \emph{single} hash by 
splitting it at a chosen bit position. This means the bucket index and 
fingerprint together always reconstruct the full hash, and expansion 
simply moves the split point.

\section{The Quotienting Idea}

\begin{definitionbox}[Quotienting]
Given a $p$-bit hash $h(x)$, we split it into two parts at bit position 
$q$:
\begin{itemize}
  \item \textbf{Quotient}: the top $q$ bits, denoted $h_q(x)$ 
        $\rightarrow$ determines the \emph{slot address}
  \item \textbf{Remainder}: the bottom $r = p - q$ bits, denoted $h_r(x)$ 
        $\rightarrow$ stored as the \emph{fingerprint}
\end{itemize}
The table has $2^q$ slots. Together, the quotient and remainder 
reconstruct the full hash: $h(x) = h_q(x) \| h_r(x)$, where $\|$ 
denotes concatenation.
\end{definitionbox}

This clean reconstruction property is fundamentally different from the cuckoo filter's approach. 
Recall that in a cuckoo filter, the bucket index comes from 
$H(x) \mod m$ and the fingerprint comes from a separate 
portion of the hash. The modulo operation is \textbf{lossy}: multiple 
distinct hash values map to the same remainder (e.g., both 
$10 \mod 8 = 2$ and $18 \mod 8 = 2$), so knowing 
$H(x) \mod m$ does not reveal what $H(x) \mod m'$ would be for 
a different table size $m'$. This lost information is precisely what 
prevents cuckoo filters from expanding. In a quotient filter, the quotient and remainder reconstruct the original value exactly, so no 
information is ever lost.

% ---- EXAMPLE: Hash split diagram ----
\begin{figure}[H]
  \centering
  \begin{tikzpicture}
    % Full hash
    \node[font=\small\ttfamily] at (-1.5, 0) {$h(x) =$};
    
    % Quotient bits
    \foreach \i/\val in {0/1, 1/0, 2/1} {
      \node[bit, fill=AccentTeal!20] (q\i) at (\i*0.6, 0) {\val};
    }
    
    % Remainder bits
    \foreach \i/\val in {3/1, 4/0, 5/1, 6/1, 7/0} {
      \node[bit, fill=AccentOrange!20] (r\i) at (\i*0.6, 0) {\val};
    }
    
    % Braces
    \draw[brace] (q0.north west) -- (q2.north east) 
      node[midway, above=10pt, font=\small, text=AccentTeal] 
      {Quotient $h_q(x)$ ($q = 3$ bits)};
    \draw[brace] (r3.north west) -- (r7.north east) 
      node[midway, above=10pt, yshift=20px, font=\small, text=AccentOrange] 
      {Remainder $h_r(x)$ ($r = 5$ bits)};
    
    % Annotations
    \node[lbl, align=center] at (1*0.6, -1.1) 
      {$\rightarrow$ slot address\\$\quot{101}_2 = 5$};
    \node[lbl, align=center, xshift=20px] at (5*0.6, -1.1) 
      {$\rightarrow$ stored fingerprint\\$\rem{10110}_2$};
  \end{tikzpicture}
  \caption{An 8-bit hash split into a 3-bit quotient (slot address) and 
           5-bit remainder (stored fingerprint). The table has 
           $2^3 = 8$ slots.}
  \label{fig:quotienting}
\end{figure}

\begin{insightbox}[Cuckoo vs.\ Quotient: The Key Difference]
In a \textbf{cuckoo filter}, the bucket index and fingerprint are 
\emph{independent}, which means knowing one tells you nothing about the other. 
In a \textbf{quotient filter}, they are \emph{complementary}, and 
together they reconstruct the full hash. This complementarity is what 
enables expansion without access to the original keys.
\end{insightbox}

\subsection{Worked Example: Insertion}

Consider a quotient filter with $p = 8$-bit hashes, $q = 3$-bit 
quotients, and $r = 5$-bit remainders. The table has $2^3 = 8$ slots. 
Suppose we insert three elements:

\begin{align*}
  h(\texttt{"hello"}) &= \underbrace{101}_{\text{quotient}} 
                         \underbrace{10101}_{\text{remainder}} 
                         \quad \rightarrow \text{slot } 5, 
                         \text{ store } \rem{10101} \\
  h(\texttt{"world"}) &= \underbrace{011}_{\text{quotient}} 
                         \underbrace{00110}_{\text{remainder}} 
                         \quad \rightarrow \text{slot } 3, 
                         \text{ store } \rem{00110} \\
  h(\texttt{"cat"})   &= \underbrace{101}_{\text{quotient}} 
                         \underbrace{11010}_{\text{remainder}} 
                         \quad \rightarrow \text{slot } 5, 
                         \text{ store } \rem{11010}
\end{align*}

The first two elements go directly into their home slots. But 
\texttt{"cat"} also hashes to slot 5, which is already occupied by 
\texttt{"hello"}. This is a \textbf{soft collision}: different keys with 
the same quotient but different remainders. The quotient filter resolves 
this by storing \texttt{"cat"}'s remainder in the next available slot 
(slot 6) using linear probing. Remainders within a run are kept in 
sorted order, so \texttt{"hello"}'s remainder (10101) stays in the 
home slot and \texttt{"cat"}'s larger remainder (11010) is placed after it:

% ---- Diagram: Simple insertion with collision ----
\begin{figure}[H]
  \centering
  \begin{tikzpicture}
    % Slots
    \foreach \i in {0,...,7} {
      \node[emptyslot, minimum width=1.3cm] (s\i) at (\i*1.4, 0) {};
      \node[font=\tiny, text=MedGray] at (\i*1.4, -0.55) {\i};
    }
    
    % Filled slots
    \node[fullslot, minimum width=1.3cm] at (3*1.4, 0) {\rem{00110}};
    \node[fullslot, minimum width=1.3cm] at (5*1.4, 0) {\rem{10101}};
    \node[fullslot, minimum width=1.3cm] at (6*1.4, 0) {\rem{11010}};
    
    % Labels
    \node[font=\scriptsize, text=AccentTeal] at (3*1.4, 0.6) 
      {\texttt{"world"}};
    \node[font=\scriptsize, text=AccentTeal] at (5*1.4, 0.6) 
      {\texttt{"hello"}};
    \node[font=\scriptsize, text=AccentTeal] at (6*1.4, 0.6) 
      {\texttt{"cat"}};
    
    % Show that cat was shifted
    \draw[arrow, AccentRed, thick] (5*1.4 - 0.1, -0.9) -- (6*1.4 + 0.1, -0.9)
      node[midway, below, font=\scriptsize, text=AccentRed, yshift=-3pt] 
      {shifted right};
    
    % Brace for run
    \draw[brace] (5*1.4 - 0.6, 1.0) -- (6*1.4 + 0.6, 1.0)
      node[midway, above=10pt, font=\scriptsize, text=DarkGray] 
      {run for quotient 5};
  \end{tikzpicture}
  \caption{After inserting \texttt{"hello"}, \texttt{"world"}, and 
           \texttt{"cat"}. Elements with the same quotient form a 
           \textbf{run} stored in contiguous slots.}
  \label{fig:quotient-insertion}
\end{figure}

The set of remainders sharing the same quotient is called a \textbf{run}. 
A sequence of contiguous runs is called a \textbf{cluster}; as the table 
fills, clusters grow longer, which increases lookup time, a limitation 
we will revisit.


\section{Metadata Bits}

Storing remainders in shifted positions creates an ambiguity: given a 
remainder in some slot, how do we know which quotient it belongs to? The 
quotient filter resolves this with three metadata bits per slot 
\cite{quotient}:

\begin{definitionbox}[Metadata Bits]
Each slot $j$ in the quotient filter maintains three single-bit flags:
\begin{enumerate}
  \item \textbf{\texttt{is\_occupied}}: Set to 1 if quotient $j$ is the 
        home slot for \emph{some} element in the filter (the element may 
        be stored elsewhere due to shifting).
  \item \textbf{\texttt{is\_continuation}}: Set to 1 if this slot holds 
        a remainder that is \emph{not the first} in its run (i.e., it 
        continues a run from a previous slot).
  \item \textbf{\texttt{is\_shifted}}: Set to 1 if the remainder in this 
        slot has been displaced from its home slot (i.e., it does not 
        occupy the slot matching its quotient).
\end{enumerate}
\end{definitionbox}

These three bits encode enough information to reconstruct which quotient 
each remainder belongs to, without storing the quotient explicitly. The 
key relationships are:

\begin{itemize}
  \item A slot with \texttt{is\_occupied = 0}, \texttt{is\_continuation = 0}, 
        \texttt{is\_shifted = 0} is \textbf{empty}.
  \item A slot with \texttt{is\_shifted = 0} and 
        \texttt{is\_continuation = 0} marks the \textbf{start of a 
        cluster}, meaning it is the first element of a run in its home slot.
  \item \texttt{is\_continuation = 0} marks the \textbf{start of a new 
        run} within a cluster.
  \item \texttt{is\_occupied} is attached to the \emph{slot index}, not 
        the stored remainder. It says ``some element with this quotient 
        exists somewhere in the table.''
\end{itemize}

\subsection{Worked Example: Metadata in Action}

Continuing from Figure~\ref{fig:quotient-insertion}, suppose we also 
insert an element with quotient 6:

\begin{align*}
  h(\texttt{"dog"}) &= \underbrace{110}_{\text{quotient}} 
                       \underbrace{01001}_{\text{remainder}} 
                       \quad \rightarrow \text{slot } 6
\end{align*}

Slot 6 is occupied by \texttt{"cat"} (which was shifted from slot 5). 
So \texttt{"dog"}'s remainder gets pushed to slot 7. The full state:

% ---- Diagram: Metadata bits ----
\begin{figure}[H]
  \centering
  \begin{tikzpicture}
    % Column headers
    \node[font=\scriptsize\bfseries, text=DarkGray] at (-2.2, 0) {Slot};
    \node[font=\scriptsize\bfseries, text=DarkGray] at (-2.2, -0.7) {Remainder};
    \node[font=\scriptsize\bfseries, text=DarkGray] at (-2.2, -1.4) 
      {\texttt{is\_occupied}};
    \node[font=\scriptsize\bfseries, text=DarkGray] at (-2.2, -2.1) 
      {\texttt{is\_cont.}};
    \node[font=\scriptsize\bfseries, text=DarkGray] at (-2.2, -2.8) 
      {\texttt{is\_shifted}};
    \node[font=\scriptsize\bfseries, text=DarkGray] at (-2.2, -3.6) 
      {Belongs to};
    
    % Slot 3 — "world"
    \node[font=\scriptsize] at (0, 0) {3};
    \node[fullslot, minimum width=1.3cm] at (0, -0.7) {\rem{00110}};
    \node[font=\scriptsize] at (0, -1.4) {1};
    \node[font=\scriptsize] at (0, -2.1) {0};
    \node[font=\scriptsize] at (0, -2.8) {0};
    \node[font=\scriptsize, text=AccentTeal] at (0, -3.6) {$q = 3$};
    
    % Slot 4 — empty
    \node[font=\scriptsize] at (1.6, 0) {4};
    \node[emptyslot, minimum width=1.3cm] at (1.6, -0.7) {};
    \node[font=\scriptsize, text=MedGray] at (1.6, -1.4) {0};
    \node[font=\scriptsize, text=MedGray] at (1.6, -2.1) {0};
    \node[font=\scriptsize, text=MedGray] at (1.6, -2.8) {0};
    \node[font=\scriptsize, text=MedGray] at (1.6, -3.6) {---};
    
    % Slot 5 — "hello"
    \node[font=\scriptsize] at (3.2, 0) {5};
    \node[fullslot, minimum width=1.3cm] at (3.2, -0.7) {\rem{10101}};
    \node[font=\scriptsize] at (3.2, -1.4) {1};
    \node[font=\scriptsize] at (3.2, -2.1) {0};
    \node[font=\scriptsize] at (3.2, -2.8) {0};
    \node[font=\scriptsize, text=AccentTeal] at (3.2, -3.6) {$q = 5$};
    
    % Slot 6 — "cat" (shifted from 5)
    \node[font=\scriptsize] at (4.8, 0) {6};
    \node[fullslot, minimum width=1.3cm] at (4.8, -0.7) {\rem{11010}};
    \node[font=\scriptsize] at (4.8, -1.4) {1};
    \node[font=\scriptsize] at (4.8, -2.1) {1};
    \node[font=\scriptsize] at (4.8, -2.8) {1};
    \node[font=\scriptsize, text=AccentTeal] at (4.8, -3.6) {$q = 5$};
    
    % Slot 7 — "dog" (shifted from 6)
    \node[font=\scriptsize] at (6.4, 0) {7};
    \node[fullslot, minimum width=1.3cm] at (6.4, -0.7) {\rem{01001}};
    \node[font=\scriptsize] at (6.4, -1.4) {0};
    \node[font=\scriptsize] at (6.4, -2.1) {0};
    \node[font=\scriptsize] at (6.4, -2.8) {1};
    \node[font=\scriptsize, text=AccentTeal] at (6.4, -3.6) {$q = 6$};
    
    % Run braces
    \draw[brace] (3.2 - 0.6, 0.4) -- (4.8 + 0.6, 0.4)
      node[midway, above=10pt, font=\scriptsize, text=AccentBlue] 
      {run for $q = 5$};
    \draw[decorate, decoration={brace, amplitude=4pt, raise=2pt, 
      mirror}, thick, DarkGray] 
      (3.2 - 0.6, -4.1) -- (6.4 + 0.6, -4.1)
      node[midway, below=8pt, font=\scriptsize, text=DarkGray] 
      {cluster (contiguous occupied region)};
  \end{tikzpicture}
  \caption{Quotient filter state after inserting \texttt{"world"} 
           ($q = 3$), \texttt{"hello"} ($q = 5$), \texttt{"cat"} 
           ($q = 5$), and \texttt{"dog"} ($q = 6$). The metadata bits 
           encode run boundaries and shift status.}
  \label{fig:quotient-metadata}
\end{figure}

Reading the metadata for each slot:
\begin{itemize}
  \item \textbf{Slot 3}: \texttt{is\_occupied = 1} (quotient 3 has 
        elements), \texttt{is\_shifted = 0} (it is in its home slot), 
        \texttt{is\_continuation = 0} (first in its run). This is 
        \texttt{"world"}'s home.
  \item \textbf{Slot 5}: \texttt{is\_occupied = 1} (quotient 5 has 
        elements), \texttt{is\_shifted = 0} (in its home slot), 
        \texttt{is\_continuation = 0} (starts a new run). This is 
        \texttt{"hello"}.
  \item \textbf{Slot 6}: \texttt{is\_occupied = 1} (quotient 6 has 
        elements somewhere), \texttt{is\_shifted = 1} (not in its home 
        slot), \texttt{is\_continuation = 1} (continues the run from 
        slot 5). This is \texttt{"cat"}, which belongs to quotient 5 but 
        is stored in slot 6.
  \item \textbf{Slot 7}: \texttt{is\_occupied = 0} (no element has 
        quotient 7), \texttt{is\_shifted = 1} (displaced from slot 6), 
        \texttt{is\_continuation = 0} (starts a new run). This is 
        \texttt{"dog"}, which belongs to quotient 6.
\end{itemize}

\subsection{Lookup Algorithm}

To query whether an element $x$ with quotient $q$ and remainder $r$ is 
in the filter:

\begin{enumerate}
  \item Check slot $q$. If \texttt{is\_occupied = 0}, then no element 
        with quotient $q$ exists, return \textsc{no}.
  \item \textbf{Find the cluster start}: scan left from slot $q$ until 
        a slot with \texttt{is\_shifted = 0} is found. This is where 
        the cluster begins.
  \item \textbf{Skip preceding runs}: scan right through the cluster, 
        counting runs. Each \texttt{is\_occupied = 1} slot to the left 
        of $q$ indicates a run to skip. Each 
        \texttt{is\_continuation = 0} slot marks the start of a new run. 
        When the count reaches zero, the current position is the start 
        of quotient $q$'s run.
  \item \textbf{Search the run}: compare the remainder in each slot of 
        the run to $r$. If a match is found, return \textsc{yes} 
        (probably in the set). If the run ends without a match, return 
        \textsc{no} (definitely not in the set).
\end{enumerate}

\begin{examplebox}[Lookup for \texttt{"dog"} (quotient $= 6$)]
\begin{enumerate}
  \item Slot 6: \texttt{is\_occupied = 1} $\rightarrow$ quotient 6 has 
        elements. Proceed.
  \item Scan left: slot 6 has \texttt{is\_shifted = 1}, slot 5 has 
        \texttt{is\_shifted = 0} $\rightarrow$ cluster starts at slot 5.
  \item Scan right from slot 5, counting runs. Between slot 5 and slot 6, 
        there is one \texttt{is\_occupied} slot (slot 5) before our 
        target (slot 6), so we skip one run:
        \begin{itemize}
          \item Slot 5: \texttt{is\_continuation = 0} $\rightarrow$ 
                start of run for $q = 5$ (skip this run)
          \item Slot 6: \texttt{is\_continuation = 1} $\rightarrow$ 
                still $q = 5$'s run
          \item Slot 7: \texttt{is\_continuation = 0} $\rightarrow$ 
                new run starts $\rightarrow$ this is $q = 6$'s run!
        \end{itemize}
  \item Slot 7 contains remainder $\rem{01001}$. Compare to 
        \texttt{"dog"}'s remainder $\rem{01001}$ $\rightarrow$ match! 
        Return \textsc{yes}.
\end{enumerate}
\end{examplebox}

Notice that this lookup required scanning through multiple slots. In the 
worst case, if all elements hash to the same quotient, a single run 
could span the entire table, requiring $O(n)$ time. In practice with a 
good hash function, runs are short and expected lookup time is $O(1)$. 
However, as the load factor increases, clusters grow and performance 
degrades; quotient filters are typically not filled beyond 75\% 
\cite{quotient}.


\section{How Expansion Works}

This is where the quotient filter's design pays off. Because the 
quotient and remainder together form the complete hash, expanding the 
table requires no access to the original keys, only a reinterpretation 
of the stored information.

\subsection{The Expansion Procedure}

To double the table from $2^q$ to $2^{q+1}$ slots:

\begin{enumerate}
  \item \textbf{Move the split point}: the quotient grows from $q$ bits 
        to $q + 1$ bits, and the remainder shrinks from $r$ bits to 
        $r - 1$ bits.
  \item \textbf{Redistribute entries}: for each stored remainder, the 
        topmost bit of the old remainder becomes part of the new quotient. 
        If this bit is 0, the entry stays in its current slot. If it is 
        1, the entry moves to the corresponding slot in the new (upper) 
        half of the table.
  \item \textbf{Truncate remainders}: remove the top bit from each 
        remainder, since it is now part of the quotient.
\end{enumerate}

\begin{examplebox}[Expansion Step by Step]
Before expansion ($q = 3$, $r = 5$, table has $2^3 = 8$ slots):
\[
  h(\texttt{"hello"}) = \underbrace{101}_{\text{quotient}} 
                         \underbrace{10101}_{\text{remainder}}
  \quad \rightarrow \text{slot } 5, \text{ store } \rem{10101}
\]

After expansion ($q = 4$, $r = 4$, table has $2^4 = 16$ slots):
\[
  h(\texttt{"hello"}) = \underbrace{1011}_{\text{new quotient}} 
                         \underbrace{0101}_{\text{new remainder}}
  \quad \rightarrow \text{slot } 11, \text{ store } \rem{0101}
\]

The hash \texttt{10110101} did not change. We simply read one more bit 
as part of the quotient. The top bit of the old remainder (\texttt{1}) 
became the new lowest bit of the quotient, changing the slot from 
$\texttt{101}_2 = 5$ to $\texttt{1011}_2 = 11$.
\end{examplebox}

% ---- Diagram: Expansion ----
\begin{figure}[H]
  \centering
  \begin{tikzpicture}
    % BEFORE: show the hash split
    \node[lbl, font=\bfseries] at (-2.5, 2.5) {Before ($q = 3$, $r = 5$):};
    
    \node[font=\small\ttfamily] at (-1.5, 1.7) {$h(x) =$};
    \foreach \i/\val in {0/1, 1/0, 2/1} {
      \node[bit, fill=AccentTeal!20] at (\i*0.5, 1.7) {\val};
    }
    \foreach \i/\val in {3/1, 4/0, 5/1, 6/0, 7/1} {
      \node[bit, fill=AccentOrange!20] at (\i*0.5, 1.7) {\val};
    }
    \draw[brace] (0*0.5-0.2, 2.15) -- (2*0.5+0.2, 2.15)
      node[midway, above=10pt, font=\scriptsize, text=AccentTeal] {$q$: slot 5};
    \draw[brace] (3*0.5-0.2, 2.15) -- (7*0.5+0.2, 2.15)
      node[midway, above=10pt, font=\scriptsize, text=AccentOrange] 
      {$r$: \rem{10101}};
    
    % Arrow
    \draw[arrow, thick, AccentBlue] (2, 1.0) -- (2, 0.3)
  node[midway, right=4pt, font=\small, align=left] {steals 1 bit\\from r to q};

    
    % AFTER: new split
    \node[lbl, font=\bfseries] at (-2.5, -0.2) {After ($q = 4$, $r = 4$):};
    
    \node[font=\small\ttfamily] at (-1.5, -1.0) {$h(x) =$};
    \foreach \i/\val in {0/1, 1/0, 2/1, 3/1} {
      \node[bit, fill=AccentTeal!20] at (\i*0.5, -1.0) {\val};
    }
    \foreach \i/\val in {4/0, 5/1, 6/0, 7/1} {
      \node[bit, fill=AccentOrange!20] at (\i*0.5, -1.0) {\val};
    }
    \draw[brace] (0*0.5-0.2, -0.55) -- (3*0.5+0.2, -0.55)
      node[midway, above=10pt, font=\scriptsize, text=AccentTeal] 
      {$q$: slot 11};
    \draw[brace] (4*0.5-0.2, -0.55) -- (7*0.5+0.2, -0.55)
      node[midway, above=10pt, font=\scriptsize, text=AccentOrange] 
      {$r$: \rem{0101}};
    
    % Highlight the stolen bit
    \draw[AccentRed, very thick, rounded corners=2pt] 
      (3*0.5-0.22, 1.7-0.28) rectangle (3*0.5+0.22, 1.7+0.28);
    \draw[AccentRed, very thick, rounded corners=2pt] 
      (3*0.5-0.22, -1.0-0.28) rectangle (3*0.5+0.22, -1.0+0.28);
    \draw[arrow, AccentRed, dashed] (3*0.5+0.3, 1.42) -- (3*0.5+0.3, -0.72);
    
    % Summary
    \node[lbl, align=left, font=\small] at (6.5, 0.5) {
      Same hash: \texttt{10110101}\\
      Slot: $5 \rightarrow 11$\\
      Remainder: 5 bits $\rightarrow$ 4 bits\\
      No original key needed!
    };
  \end{tikzpicture}
  \caption{Expansion moves the split point one bit to the right. The top 
           bit of the old remainder (highlighted) becomes the new lowest 
           bit of the quotient. The hash itself is unchanged.}
  \label{fig:expansion-split}
\end{figure}

% ---- Diagram: Table before and after ----
\begin{figure}[H]
  \centering
  \begin{tikzpicture}[node distance=0pt]
    % BEFORE expansion
    \node[lbl, font=\bfseries] at (1.5, 1.8) {Before Expansion (8 slots)};
    
    \foreach \i/\fp in {0/, 1/, 2/, 3/\rem{00110}, 4/, 
                        5/\rem{10101}, 6/\rem{11010}, 7/} {
      \ifx\fp\empty
        \node[emptyslot, minimum width=1.3cm] (a\i) at (\i*1.4, 1) {};
      \else
        \node[fullslot, minimum width=1.3cm] (a\i) at (\i*1.4, 1) {\fp};
      \fi
      \node[font=\tiny, text=MedGray] at (\i*1.4, 0.45) {\i};
    }
    
    % Arrow down
\begin{scope}[yshift=-10pt]
  \draw[arrow, thick, AccentBlue] (4.9, 0.5) -- (4.9, -0.5) 
    node[midway, right, font=\small] 
    {expand: steal 1 bit from each remainder};
\end{scope}    
    % AFTER expansion (16 slots, show only relevant ones)
  \begin{scope}[yshift=-10pt]
     \node[lbl, font=\bfseries] at (4.5, -1.0) {After Expansion (16 slots)};
    
    \foreach \i/\fp in {0/, 1/, 2/, 3/, 4/, 5/, 6/\rem{0110}, 
                        7/, 8/, 9/, 10/, 11/\rem{0101}, 
                        12/, 13/\rem{1010}} {
      \ifx\fp\empty
        \node[emptyslot, minimum width=0.75cm] (b\i) at (\i*0.8, -1.8) {};
      \else
        \node[fullslot, minimum width=0.75cm] (b\i) at (\i*0.8, -1.8) {\fp};
      \fi
      \node[font=\tiny, text=MedGray] at (\i*0.8, -2.35) {\i};
    }
    
    \node[lbl, align=center] at (4.5, -3.1) {
      Remainders shrink: 5 bits $\rightarrow$ 4 bits.
      Entries redistribute across doubled table.
    };
  \end{scope}
  \end{tikzpicture}
  \caption{Quotient filter expansion: the table doubles from 8 to 16 
           slots. Each remainder loses its top bit to the quotient.
           \texttt{"hello"} moves from slot 5 to slot 11; 
           \texttt{"cat"} moves from slot 6 to slot 13.}
  \label{fig:expansion}
\end{figure}

\begin{insightbox}[Why This Works Without Original Keys]
The expansion procedure never needs the original element $x$ or the 
full hash $\text{hash}(x)$. It only needs the \emph{quotient} (known 
from the slot position) and the \emph{remainder} (stored in the slot). 
Together these reconstruct the full $p$-bit hash, which can then be 
re-split at the new position. This is the fundamental advantage of 
quotienting over the cuckoo filter's modulo-based addressing.
\end{insightbox}


\section{The Void Entry Problem}

Expansion comes at a cost: each time the table doubles, every remainder 
loses one bit. For an entry that was present at the very first expansion 
and survives through $X$ expansions, its remainder shrinks from $r$ bits 
to $r - X$ bits.

\subsection{How Void Entries Arise}

Consider an entry inserted with an initial remainder of $r = 5$ bits. 
After each expansion, the top bit of its remainder is consumed by the 
growing quotient:

% ---- Diagram: Void entry formation ----
\begin{figure}[H]
  \centering
  \begin{tikzpicture}
    \node[lbl, font=\bfseries] at (-2.5, 1) {Expansion};
    \node[lbl, font=\bfseries] at (1.5, 1) {Remainder};
    \node[lbl, font=\bfseries] at (4.5, 1) {Bits left};
    
    % Row 0
    \node[lbl] at (-2.5, 0) {0 (initial)};
    \node[bit, fill=BitOne] at (0.5, 0) {1};
    \node[bit, fill=BitOne] at (1.0, 0) {0};
    \node[bit, fill=BitOne] at (1.5, 0) {1};
    \node[bit, fill=BitOne] at (2.0, 0) {0};
    \node[bit, fill=BitOne] at (2.5, 0) {1};
    \node[lbl] at (4.5, 0) {5};
    
    % Row 1
    \node[lbl] at (-2.5, -0.8) {1};
    \node[bit, fill=BitOne] at (1.0, -0.8) {0};
    \node[bit, fill=BitOne] at (1.5, -0.8) {1};
    \node[bit, fill=BitOne] at (2.0, -0.8) {0};
    \node[bit, fill=BitOne] at (2.5, -0.8) {1};
    \node[bit, fill=BitZero, opacity=0.3] at (0.5, -0.8) {};
    \node[lbl] at (4.5, -0.8) {4};
    
    % Row 2
    \node[lbl] at (-2.5, -1.6) {2};
    \node[bit, fill=BitOne] at (1.5, -1.6) {1};
    \node[bit, fill=BitOne] at (2.0, -1.6) {0};
    \node[bit, fill=BitOne] at (2.5, -1.6) {1};
    \node[bit, fill=BitZero, opacity=0.3] at (0.5, -1.6) {};
    \node[bit, fill=BitZero, opacity=0.3] at (1.0, -1.6) {};
    \node[lbl] at (4.5, -1.6) {3};
    
    % Row 3
    \node[lbl] at (-2.5, -2.4) {3};
    \node[bit, fill=BitOne] at (2.0, -2.4) {0};
    \node[bit, fill=BitOne] at (2.5, -2.4) {1};
    \node[bit, fill=BitZero, opacity=0.3] at (0.5, -2.4) {};
    \node[bit, fill=BitZero, opacity=0.3] at (1.0, -2.4) {};
    \node[bit, fill=BitZero, opacity=0.3] at (1.5, -2.4) {};
    \node[lbl] at (4.5, -2.4) {2};
    
    % Row 4
    \node[lbl] at (-2.5, -3.2) {4};
    \node[bit, fill=BitOne] at (2.5, -3.2) {1};
    \node[bit, fill=BitZero, opacity=0.3] at (0.5, -3.2) {};
    \node[bit, fill=BitZero, opacity=0.3] at (1.0, -3.2) {};
    \node[bit, fill=BitZero, opacity=0.3] at (1.5, -3.2) {};
    \node[bit, fill=BitZero, opacity=0.3] at (2.0, -3.2) {};
    \node[lbl] at (4.5, -3.2) {1};
    
    % Row 5 — VOID
    \node[lbl, text=AccentRed, font=\bfseries] at (-2.5, -4.0) {5};
    \node[bit, fill=VoidColor] at (0.5, -4.0) {?};
    \node[bit, fill=VoidColor] at (1.0, -4.0) {?};
    \node[bit, fill=VoidColor] at (1.5, -4.0) {?};
    \node[bit, fill=VoidColor] at (2.0, -4.0) {?};
    \node[bit, fill=VoidColor] at (2.5, -4.0) {?};
    \node[lbl, text=AccentRed, font=\bfseries] at (4.5, -4.0) {0 = \void{}};
    
    % Arrow
    \draw[arrow, AccentRed, very thick] (5.5, -4.0) -- (7.5, -4.0) 
      node[right, font=\small, text=AccentRed, align=left] {
        No fingerprint bits left!\\
        Matches \textbf{any} query\\
        $\rightarrow$ 100\% false positive};
  \end{tikzpicture}
  \caption{A remainder shrinks by one bit per expansion. After $r$ 
           expansions, zero bits remain, the entry becomes \void{}.}
  \label{fig:void-formation}
\end{figure}

\begin{definitionbox}[Void Entry]
A \textbf{void entry} is an entry whose remainder has been reduced to 
zero bits through repeated expansions. Since there is no fingerprint 
left to compare against, a void entry matches \emph{every} query to 
its slot, producing a guaranteed false positive. Formally, if an entry 
was created with remainder length $r$ and has survived $X \geq r$ 
expansions, its current remainder length is $\max(0, r - X) = 0$.
\end{definitionbox}

\subsection{The Expansion Dilemma}

Void entries create a specific problem during expansion. When the table 
doubles, each entry moves to one of two possible slots in the new table, determined by the top bit of its current remainder. But a void entry 
has \emph{no remaining bits}. Without a bit to read, the filter cannot 
determine which of the two new slots the entry belongs in:

% ---- Diagram: Void entry expansion problem ----
\begin{figure}[H]
  \centering
  \begin{tikzpicture}
    % Before
    \node[lbl, font=\bfseries] at (1.5, 2) {Before expansion (8 slots)};
    \foreach \i in {0,...,7} {
      \node[emptyslot, minimum width=0.9cm] (a\i) at (\i*1.0, 1.2) {};
      \node[font=\tiny, text=MedGray] at (\i*1.0, 0.75) {\i};
    }
    \node[voidslot, minimum width=0.9cm] at (3*1.0, 1.2) {V};
    
    % After
    \node[lbl, font=\bfseries] at (5, -2.7) {After expansion (16 slots)};
    \foreach \i in {0,...,15} {
      \node[emptyslot, minimum width=0.55cm] (b\i) at (\i*0.6, -1.5) {};
      \node[font=\tiny, text=MedGray] at (\i*0.6, -1.95) {\i};
    }
    
    % Two possible destinations
    \node[voidslot, minimum width=0.55cm] at (6*0.6, -1.5) {?};
    \node[voidslot, minimum width=0.55cm] at (7*0.6, -1.5) {?};
    
    % Arrows to both
    \draw[dasharrow, AccentRed] (3*1.0, 0.85) -- (6*0.6, -1.15)
      node[midway, left, xshift=-4pt, font=\scriptsize, text=AccentRed] {bit = 0?};
    \draw[dasharrow, AccentRed] (3*1.0, 0.85) -- (7*0.6, -1.15)
      node[midway, right, xshift=4pt, font=\scriptsize, text=AccentRed] {bit = 1?};
    
    % Annotation
    \node[lbl, text=AccentRed, font=\small, align=center] at (12, -0.1) 
      {No remainder bits left\\to decide which slot!\\\\
       Slot 6 (if bit was 0)\\or slot 7 (if bit was 1)\\
       \textbf{We don't know.}};
  \end{tikzpicture}
  \caption{A void entry at slot 3 must move to either slot 6 or slot 7 
           during expansion, but with zero remainder bits, there is no 
           information to determine which. This is the core problem that 
           InfiniFilter and the Aleph Filter each solve differently.}
  \label{fig:void-expansion}
\end{figure}

\subsection{InfiniFilter's Approach: Secondary Tables}

InfiniFilter \cite{infinifilter} solves this by placing void entries into 
\textbf{secondary hash tables}, one per expansion level. When an entry 
becomes void, it is moved to a secondary table that stores enough 
additional hash bits to track its position.

However, this means that queries for non-existent or old keys must 
now search the main table \emph{and} all secondary tables:

% ---- Diagram: InfiniFilter multi-table query ----
\begin{figure}[H]
  \centering
  \begin{tikzpicture}
    % Main table
    \node[lbl, font=\bfseries] at (-1.5, 0) {Main table};
    \foreach \i in {0,...,7} {
      \node[emptyslot, minimum width=0.7cm] (m\i) at (\i*0.8, -0.7) {};
    }
    
    % Secondary table 1
    \node[lbl, font=\bfseries] at (-1.5, -1.8) {Secondary 1};
    \foreach \i in {0,...,3} {
      \node[emptyslot, minimum width=0.7cm, fill=orange!8] 
        (s\i) at (\i*0.8, -2.5) {};
    }
    
    % Secondary table 2
    \node[lbl, font=\bfseries] at (-1.5, -3.6) {Secondary 2};
    \foreach \i in {0,...,1} {
      \node[emptyslot, minimum width=0.7cm, fill=red!8] 
        (t\i) at (\i*0.8, -4.3) {};
    }
    
    % Query path
    \node[hashbox] (q) at (8, 0) {query $x$};
    \draw[arrow, AccentBlue] (q.west) -- (m7.east)
      node[midway, above, yshift=10pt, font=\scriptsize] {check};
    \draw[dasharrow, AccentRed] (m7.south) -- (s3.north east)
      node[midway, yshift=-6pt, right, font=\scriptsize, text=AccentRed] {miss};
    \draw[dasharrow, AccentRed] (s3.south) -- (t1.north east)
      node[midway, yshift=-6pt, right, font=\scriptsize, text=AccentRed] {miss};
    
    % Annotation
    \node[lbl, text=AccentRed, align=left, font=\small] at (7.5, -3.5) {
      Must check \textbf{every} table\\
      Query time: $\bigO{\log(N)/F}$\\
      (one table per expansion)
    };
  \end{tikzpicture}
  \caption{InfiniFilter stores void entries in secondary hash tables. 
           In the fixed-width regime, queries may need to search all tables,
           degrading to $\bigO{\log(N)/F}$ (equivalently $\bigO{\log N}$ for
           constant $F$).}
  \label{fig:infinifilter}
\end{figure}

With each expansion adding a new secondary table, the number of tables 
grows as $\bigO{\log(N)/F}$ in the fixed-width regime (where $N$ is the 
total number of insertions and $F$ is the base fingerprint length). 
A query that targets a non-existent key or an old key must search every 
table, giving worst-case $\bigO{\log(N)/F}$ query time (or $\bigO{\log N}$ 
for constant $F$), a significant 
degradation from the $\bigO{1}$ queries of the original quotient filter.

\subsection{The Aleph Filter's Solution: Void Duplication}

The Aleph Filter \cite{aleph} takes a fundamentally different approach. 
Instead of moving void entries to secondary tables, it \textbf{duplicates} 
each void entry into \emph{both} possible slots during expansion. The two candidate slots in the expanded table are $2q$ and $2q + 1$, 
corresponding to appending a $0$ or $1$ bit to the old quotient:
\[
  q \| 0 = 2q \qquad \text{and} \qquad q \| 1 = 2q + 1
\]
where $\|$ denotes bit concatenation. A comparison betwen InfiniFilter and Aleph Filter:

% ---- Diagram: Aleph vs InfiniFilter ----
\begin{figure}[H]
  \centering
  \begin{tikzpicture}[scale=0.9, every node/.style={transform shape}]
    % BEFORE
    \node[lbl, font=\bfseries] at (1.6, 1.5) {Before Expansion};
    \foreach \i/\content/\style in {0//emptyslot, 1//emptyslot, 
                                     2//emptyslot, 3/V/voidslot} {
      \node[\style, minimum width=1cm] (a\i) at (\i*1.2, 0.7) {\content};
      \node[font=\tiny, text=MedGray] at (\i*1.2, 0.15) {\i};
    }
    
    % Arrow
    \draw[arrow, very thick, AccentBlue] (2, -0.3) -- (2, -1.0);
    
    % AFTER — InfiniFilter
    \node[lbl, font=\bfseries] at (-3, -1.6) {InfiniFilter:};
    \foreach \i in {0,...,7} {
      \node[emptyslot, minimum width=0.8cm] (inf\i) at ({-5+\i*0.9}, -2.3) {};
    }
    \node[font=\tiny, text=MedGray] at (-5, -2.8) {0};
    \node[font=\tiny, text=MedGray] at ({-5+7*0.9}, -2.8) {7};
    
    \node[voidslot, minimum width=0.8cm] (sec) at (-3, -3.5) {V};
    \node[lbl, font=\scriptsize] at (-3, -4.1) {secondary table};
    \draw[dasharrow, AccentRed] (inf3.south) -- (sec.north);
    \node[lbl, text=AccentRed, font=\scriptsize, anchor=west] at (0, -3.5) 
      {Query: $\bigO{\log N}$};
    
    % AFTER — Aleph
    \node[lbl, font=\bfseries] at (5.5, -1.6) {Aleph Filter:};
    \foreach \i in {0,...,7} {
      \node[emptyslot, minimum width=0.8cm] (al\i) at ({3+\i*0.9}, -2.3) {};
    }
    \node[font=\tiny, text=MedGray] at (3, -2.8) {0};
    \node[font=\tiny, text=MedGray] at ({3+7*0.9}, -2.8) {7};
    
    \node[voidslot, minimum width=0.8cm] at ({3+6*0.9}, -2.3) {V};
    \node[voidslot, minimum width=0.8cm] at ({3+7*0.9}, -2.3) {V};
    
    \draw[brace] ({3+6*0.9-0.35}, -1.7) -- ({3+7*0.9+0.35}, -1.7)
      node[midway, above=10pt, font=\scriptsize, text=AccentGreen] 
      {duplicated!};
    \node[lbl, text=AccentGreen, font=\scriptsize, anchor=west, xshift=10pt] at (9.5, -2.3) 
      {Query: $\bigO{1}$};
  \end{tikzpicture}
  \caption{Aleph Filter duplicates void entries to both possible slots, 
           keeping queries $\bigO{1}$.}
  \label{fig:aleph-vs-infini}
\end{figure}

Since a void entry matches any query to its slot (it has no fingerprint 
to compare), placing a copy in both candidate slots ensures correctness: 
regardless of which slot the original element ``should'' have gone to, 
a copy is there. All entries remain in a single table, so queries never 
need to search secondary structures.

This is the core innovation of the Aleph Filter, which the next chapter 
develops in full detail. The complete design addresses three challenges:
\begin{enumerate}
  \item \textbf{Void duplication:} keeps queries $\bigO{1}$ by 
        eliminating secondary tables.
  \item \textbf{Tombstone-based deletion:} handles the complication that 
        deleting a duplicated void entry must not remove both copies.
  \item \textbf{Rejuvenation:} prevents void entries from accumulating 
        indefinitely by periodically replacing them with fresh 
        fingerprints.
\end{enumerate}

\begin{warningbox}[The Complete Filter Evolution]
\begin{center}
\begin{tabular}{l c c c c}
  \hline
  & \textbf{Delete} & \textbf{Expand} & \textbf{Query} & \textbf{Stable FPR} \\
  \hline
  Bloom filter     & \texttimes & \texttimes & $\bigO{k}$       & \texttimes \\
  Cuckoo filter    & \checkmark & \texttimes & $\bigO{1}$       & \texttimes \\
  Quotient filter  & \checkmark & \checkmark & $\bigO{1}$*      & \texttimes \\
  InfiniFilter     & \checkmark & \checkmark & $\bigO{\log(N)/F}$  & \checkmark \\
  \textbf{Aleph Filter} & \checkmark & \checkmark & $\bigO{1}$ & \checkmark \\
  \hline
\end{tabular}

\smallskip
{\small *amortized; degrades at high load factors}
\end{center}
Each row solves a limitation of the row above it. The Aleph Filter is 
the first to achieve all four properties simultaneously.
\end{warningbox}


% ============================================================================
% CHAPTER 3: THE ALEPH FILTER
% ============================================================================
